\documentclass[letterpaper]{article}
\usepackage[submission]{aaai2027}
\usepackage{booktabs}
\usepackage{amsmath}
\usepackage{graphicx}
\usepackage{url}
\usepackage{float}

\setlength{\pdfpagewidth}{8.5in}
\setlength{\pdfpageheight}{11in}
\frenchspacing

\title{Supplementary Material for\\
InTraScale: Accelerating Video Generation via In-Trajectory Resolution Scaling}
\author{Anonymous Submission}
\affiliations{}

\begin{document}
\maketitle

\section{Scope and Evidence Conventions}

This supplement provides implementation details, protocol definitions,
additional results, and paired statistical analyses for InTraScale.  The main
paper remains self-contained.  All comparisons within a suite use matched text
prompts, generation seeds, initial noise, scheduler, guidance, frame count, and
output resolution.

\paragraph{Metric naming.}
We evaluate Subject Consistency, Background Consistency, Motion Smoothness,
Aesthetic Quality, and Imaging Quality using the VBench custom-input protocol.
We call their unweighted mean \emph{VBench-5}.  VBench-5 is a compact
paper-specific summary and is not the official VBench Quality, Semantic, or
Total score.  Paired analyses use prompt--seed pairs as statistical units.

\paragraph{Method naming and provenance.}
The immutable experiment archives predate the final terminology.  Internal
case identifiers containing \texttt{talh}, \texttt{cll}, or \texttt{stage2}
denote InTraScale, ITU-only, and ITU, respectively.  The code-and-data package
contains an explicit mapping and retains the internal identifiers so that
archived results remain traceable.

\paragraph{Evidence coverage.}
The package includes validated Wan50 and Distill4 quality, timing, component,
and paired-statistics exports.  It also includes the checksum-verified
50-sample raw \(368{\times}640\rightarrow720{\times}1248\) ITU operator
evaluation and a deterministic aggregation script.  The local archive does not
contain the complete custom ITU/TTD checkpoints or exact hardware/software
version metadata.  These missing items are not reconstructed from filenames or
earlier prose.

\section{Implementation and Training Details}

\subsection{In-Trajectory Upsampler}

ITU maps a 16-channel clean video latent to the target spatial grid without
changing the temporal dimension.  A 3D input projection and four
spatiotemporal residual blocks precede the spatial lifting operator; four
further residual blocks and a 3D output projection reconstruct the target
latent.  The hidden width is 256.  For the \(3/2\) grid ratio, a \(3\times\)
PixelShuffle is followed by anti-aliased stride-2 downsampling.  For the
approximately \(2\times\) ratio, a \(2\times\) PixelShuffle produces a
\(92\times160\) latent grid that is center-cropped to \(90\times156\).

ITU minimizes
\[
  \mathcal{L}_{\mathrm{ITU}}
  =\mathcal{L}_{\mathrm{rec}}
   +0.2\mathcal{L}_{\mathrm{lf}}
   +0.1\mathcal{L}_{\mathrm{temp}},
\]
where the reconstruction term is Charbonnier loss with
\(\epsilon=10^{-3}\), the low-frequency term preserves downsampled content,
and the temporal term matches first-order latent differences.

\subsection{Trajectory-Tail Distillation}

TTD is a transition-local LoRA update.  It targets the attention
\(q,k,v,o\) projections and the two feed-forward projections, uses rank and
training scale \(r=\alpha=32\), and keeps the base denoiser frozen.  Its loss is
\[
  \mathcal{L}_{\mathrm{TTD}}
  =\mathcal{L}_{1}+0.1\mathcal{L}_{\mathrm{MSE}}
   +\lambda_{\mathrm{tt}}\mathcal{L}_{\mathrm{temp}}.
\]
We use \(\lambda_{\mathrm{tt}}=0.05\) for Wan50 step 40 and zero for Wan50
step 45 and Distill4 step 3.  The inference strength is selected from
\(\{0.25,0.5,0.75,1.0\}\) on held-out development prompts; \(0.75\) is then
fixed for the test sets.

\begin{table}[H]
  \centering
  \small
  \setlength{\tabcolsep}{4pt}
  \begin{tabular}{lcc}
    \toprule
    Setting & ITU & TTD \\
    \midrule
    Optimizer & AdamW & AdamW \\
    Maximum steps & 50,000 & 10,000 \\
    Precision & bf16 & bf16 \\
    Learning rate & \(1{\times}10^{-4}\) & \(5{\times}10^{-5}\) \\
    Weight decay & 0.01 & 0.01 \\
    Effective batch size & 8 & 8 \\
    Gradient clipping & 1.0 & 1.0 \\
    Training seed & 1234 & 1234 \\
    EMA & 0.9999 & none \\
    Validation fraction & 0.05 & 0.02 \\
    Trainable parameters & ITU only & rank-32 LoRA only \\
    \bottomrule
  \end{tabular}
  \caption{Final training settings encoded in the supplied configurations.
  ITU uses batch size 1 with distributed ranks/accumulation totaling an
  effective batch size of 8.}
  \label{tab:training}
\end{table}

\section{Evaluation Protocol}

\paragraph{Generation.}
Wan50 uses Wan2.1 T2V-1.3B with 50 denoising evaluations and transition steps
40 or 45 for InTraScale.  Distill4 uses Wan2.1 T2V-14B
StepDistill-CfgDistill with scheduler steps
\([1000,750,500,250]\); InTraScale-D4 performs three low-resolution
evaluations and one target-resolution evaluation.  Every output contains 81
frames at \(720\times1248\).

\paragraph{Prompts and randomness.}
Quality uses ten matched prompt--seed pairs in each suite.  Wan50 uses seeds
9700--9709 and Distill4 uses seeds 9800--9809.  The state-reconstruction and
strength selection study uses eight disjoint prompts with seed base 16000.
Prompt lists are included in the code-and-data package.

\paragraph{Warm latency.}
The model and checkpoints remain resident in one process.  We run one untimed
warm-up and then five matched generations.  CUDA synchronization surrounds
each measurement.  Pipeline latency covers denoising, TTD, ITU, scheduler
state reconstruction, VAE decoding, and output writing; it excludes runner
construction and checkpoint loading.  Exact GPU, driver, CUDA, and framework
versions were not preserved in the local export and must be recovered from the
measurement machine before final submission.

\paragraph{Statistics.}
For quality, we report prompt-paired means, nonparametric bootstrap 95\%
confidence intervals, win/loss counts, and two-sided sign tests.  Timing
variation is the sample standard deviation across the five warm repeats.
These two forms of uncertainty are not interchanged.

\section{Code, Data, and Media Inventory}

The accompanying code-and-data ZIP contains the ITU model, TTD LoRA utilities,
Wan50 and Distill4 bridges, data builders, training entry points, inference
runners, VBench aggregation, paired-statistics scripts, final configurations,
prompt lists, sanitized raw metric JSONs, and the compact CSV tables used here.
The media ZIP contains two matched prompt--seed groups.  Each group compares
Trilinear, ITU-only, and InTraScale at 1248-by-720 resolution, 81 frames, and
16 frames per second.  Legacy 640-by-368 files misleadingly named
\texttt{Native-HR} are explicitly excluded.

\paragraph{Reproducibility boundary.}
Public base-model weights, large latent LMDBs, and generated-video test sets are
not duplicated.  More importantly, the current local archive does not contain
the final custom ITU/TTD weights.  The package therefore supports independent
audit of implementation, protocols, raw metric exports, and statistical
claims, but full regeneration requires adding the final custom checkpoints.

\clearpage
\onecolumn

\section{Additional Wan50 Results}

\begin{table}[H]
  \centering
  \small
  \setlength{\tabcolsep}{5pt}
  \begin{tabular}{lrrrr}
    \toprule
    Method & LR evals & HR evals & VBench-5 \(\uparrow\) &
    Warm latency (s) \(\downarrow\) \\
    \midrule
    Native-HR & 0 & 50 & 0.82836 & \(187.54\pm0.26\) \\
    Trilinear@40 & 40 & 10 & 0.77812 & \(58.12\pm0.20\) \\
    InTraScale@40 & 40 & 10 & 0.80983 & \(58.49\pm0.16\) \\
    Trilinear@45 & 45 & 5 & 0.76776 & \(41.98\pm0.19\) \\
    RALU-style@45 & 45 & 5 & 0.80440 & \(42.01\pm0.19\) \\
    InTraScale@45 & 45 & 5 & 0.80792 & \(42.26\pm0.20\) \\
    Trilinear@48 & 48 & 2 & 0.76409 & \(32.27\pm0.18\) \\
    ITU-0HR & 50 & 0 & 0.79647 & \(25.88\pm0.11\) \\
    ITU-1HR & 50 & 1 & 0.80093 & \(29.46\pm0.22\) \\
    ITU-2HR & 50 & 2 & 0.80355 & \(33.23\pm0.10\) \\
    ITU-5HR (LUVE-style) & 50 & 5 & 0.80073 & \(44.00\pm0.27\) \\
    \bottomrule
  \end{tabular}
  \caption{Wan50 warm-model operating points.  Endpoint-first ITU baselines
  complete all 50 low-resolution evaluations before lifting and optional
  target-resolution refinement.}
  \label{tab:wan50-full}
\end{table}

\subsection{ITU Cross-Resolution Reconstruction}

\begin{table}[H]
  \centering
  \scriptsize
  \setlength{\tabcolsep}{3.2pt}
  \begin{tabular}{llrrr}
    \toprule
    Input & Metric & Trilinear & ITU & Wins \\
    \midrule
    \(480p\) & Latent \(L_1\) \(\downarrow\) & 0.19288 & 0.09913 & 49/50 \\
    \(480p\) & PSNR \(\uparrow\) & 24.2927 & 29.9980 & 50/50 \\
    \(480p\) & SSIM \(\uparrow\) & 0.71588 & 0.86961 & 49/50 \\
    \(480p\) & LPIPS \(\downarrow\) & 0.23575 & 0.06294 & 50/50 \\
    \(480p\) & Temporal \(L_1\) \(\downarrow\) & 0.02121 & 0.01468 & 50/50 \\
    \(480p\) & HF-energy error \(\downarrow\) & 0.00449 & 0.00161 & 48/50 \\
    \midrule
    \(368p\) & Latent \(L_1\) \(\downarrow\) & 0.24437 & 0.16284 & 50/50 \\
    \(368p\) & PSNR \(\uparrow\) & 23.4238 & 24.1021 & 37/50 \\
    \(368p\) & SSIM \(\uparrow\) & 0.67404 & 0.75920 & 48/50 \\
    \(368p\) & LPIPS \(\downarrow\) & 0.33522 & 0.16403 & 50/50 \\
    \(368p\) & Temporal \(L_1\) \(\downarrow\) & 0.02388 & 0.02017 & 49/50 \\
    \(368p\) & HF-energy error \(\downarrow\) & 0.00792 & 0.00233 & 50/50 \\
    \bottomrule
  \end{tabular}
  \caption{ITU operator evaluation on 50 clean-latent pairs per setting.
  ``Wins'' counts samples on which ITU is better.  The raw 368p records and
  deterministic aggregation script accompany the submission.}
  \label{tab:itu-operator}
\end{table}

\subsection{TTD Endpoint Alignment}

\begin{table}[H]
  \centering
  \small
  \setlength{\tabcolsep}{4pt}
  \begin{tabular}{lrrrrr}
    \toprule
    Transition & Base \(L_1\) & TTD \(L_1\) & Reduction &
    95\% CI of improvement & Wins \\
    \midrule
    Wan50 \(s=40\) & 0.03215 & 0.02385 & 25.82\% &
    [0.00479, 0.01302] & 10/10 \\
    Wan50 \(s=45\) & 0.02363 & 0.01866 & 21.03\% &
    [0.00260, 0.00840] & 10/10 \\
    Distill4 \(s=3\) & 0.04286 & 0.04070 & 5.03\% &
    [0.00148, 0.00299] & 10/10 \\
    \bottomrule
  \end{tabular}
  \caption{Prompt-paired endpoint alignment at the selected TTD strength
  \(\lambda=0.75\).  Positive confidence intervals indicate lower endpoint
  \(L_1\) after TTD.}
  \label{tab:ttd-endpoint}
\end{table}

\subsection{Factorial Separation of ITU and TTD}

\begin{table}[H]
  \centering
  \small
  \setlength{\tabcolsep}{4pt}
  \begin{tabular}{llrrr}
    \toprule
    Suite & Effect & Before & After & \(\Delta\) VBench-5 \\
    \midrule
    Wan50@40 & ITU within matched transition & 0.77812 & 0.80896 & +0.03085 \\
    Wan50@40 & TTD within ITU & 0.80896 & 0.80983 & +0.00087 \\
    Wan50@45 & ITU within matched transition & 0.76776 & 0.80842 & +0.04066 \\
    Wan50@45 & TTD within ITU & 0.80842 & 0.80792 & \(-0.00050\) \\
    Distill4@3 & ITU within matched transition & 0.81796 & 0.85670 & +0.03874 \\
    Distill4@3 & TTD within ITU & 0.85670 & 0.85680 & +0.00010 \\
    \bottomrule
  \end{tabular}
  \caption{Factorial quality effects.  ITU supplies the dominant aggregate
  perceptual gain; TTD targets transition endpoint mismatch rather than acting
  as a general perceptual enhancer.}
  \label{tab:factorial}
\end{table}

\section{Additional Distill4 Results}

\begin{table}[H]
  \centering
  \scriptsize
  \setlength{\tabcolsep}{3.2pt}
  \begin{tabular}{lrrrrrrrr}
    \toprule
    Method & LR/HR & Subject & Background & Motion & Aesthetic & Imaging &
    VBench-5 & Latency (s) \\
    \midrule
    Native-HR4 & 0/4 & 0.97186 & 0.97557 & 0.99123 & 0.65203 & 0.71135 &
    0.86041 & 42.49 \\
    Trilinear@3 & 3/1 & 0.95932 & 0.96937 & 0.98376 & 0.58954 & 0.58783 &
    0.81796 & 18.33 \\
    InTraScale-D4@3 & 3/1 & 0.96226 & 0.97026 & 0.98815 & 0.64914 & 0.71418 &
    0.85680 & 19.92 \\
    ITU-2HR (LUVE-style) & 4/2 & 0.96238 & 0.97341 & 0.98683 & 0.65191 &
    0.72450 & 0.85981 & 28.76 \\
    MrFlow-style & 4/1 & 0.96379 & 0.97259 & 0.98796 & 0.65491 & 0.71591 &
    0.85903 & 27.12 \\
    \bottomrule
  \end{tabular}
  \caption{Distill4 warm-model quality--efficiency.  The last two methods are
  endpoint-first baselines and perform additional processing after completing
  the four-step low-resolution trajectory.}
  \label{tab:distill4-full}
\end{table}

\begin{table}[H]
  \centering
  \small
  \setlength{\tabcolsep}{4pt}
  \begin{tabular}{llrrrr}
    \toprule
    Baseline & Metric & Baseline--InTraScale & 95\% CI & Wins/Losses & \(p\) \\
    \midrule
    Trilinear@3 & VBench-5 & \(-0.03884\) & [\(-0.04840,-0.02921\)] &
    0/10 & 0.00195 \\
    Trilinear@3 & Temporal & \(-0.00264\) & [\(-0.00434,-0.00130\)] &
    1/9 & 0.02148 \\
    ITU-2HR & VBench-5 & \(+0.00301\) & [0.00077, 0.00564] &
    7/3 & 0.34375 \\
    ITU-2HR & Temporal & \(-0.00333\) & [\(-0.00547,-0.00152\)] &
    1/9 & 0.02148 \\
    MrFlow-style & VBench-5 & \(+0.00223\) & [\(-0.00277,0.00752\)] &
    6/4 & 0.75391 \\
    MrFlow-style & Temporal & \(-0.00180\) & [\(-0.00289,-0.00075\)] &
    2/8 & 0.10938 \\
    \bottomrule
  \end{tabular}
  \caption{Distill4 prompt-paired statistics.  Differences are baseline minus
  InTraScale; ``Temporal'' is temporal consistency, so higher is better.
  \(p\) is the two-sided sign-test value.}
  \label{tab:distill4-stats}
\end{table}

\subsection{Scheduler-State Reconstruction and Strength Selection}

\begin{table}[H]
  \centering
  \small
  \setlength{\tabcolsep}{4pt}
  \begin{tabular}{lrrrr}
    \toprule
    Target-grid noise & TTD strength & VBench-5 & Temporal consistency & Selected \\
    \midrule
    Seed-derived Gaussian & 0.25 & 0.83816 & 0.97476 & no \\
    Seed-derived Gaussian & 0.50 & 0.83893 & 0.97468 & no \\
    Seed-derived Gaussian & 0.75 & 0.84009 & 0.97465 & yes \\
    Seed-derived Gaussian & 1.00 & 0.83933 & 0.97464 & no \\
    Resized LR flow & 0.25 & 0.81467 & 0.95862 & no \\
    Resized LR flow & 0.50 & 0.81488 & 0.95853 & no \\
    Resized LR flow & 0.75 & 0.81526 & 0.95860 & no \\
    Resized LR flow & 1.00 & 0.81486 & 0.95859 & no \\
    \bottomrule
  \end{tabular}
  \caption{Eight-prompt development study.  Scheduler-consistent
  target-resolution noise improves both summary quality and temporal
  consistency across all tested TTD strengths.}
  \label{tab:state-reconstruction}
\end{table}

\end{document}
